\mxfront{stmg-bank-a}

% -----------------------------------------------------------------------------
% Chapitre 1 -- Suites numériques
% -----------------------------------------------------------------------------
\mxexdef{1}{Reconnaître une suite arithmétique par différence}{%
Le responsable d'un entrepôt teste trois modèles. Pour tout entier naturel $n$, $n$ désigne le numéro de la semaine.
\begin{itemize}
\item Nombre de colis préparés : $a_n=3(2n+7)-n$.
\item Nombre de réclamations clients : $b_n=2n^2-n+4$.
\item Nombre de palettes encore en stock : $c_n=150-7n$.
\end{itemize}
\begin{enumerate}[label=\alph*.]
\item Pour chacune des trois suites, calculer $u_{n+1}-u_n$ (avec $u=a$, $b$ ou $c$) et préciser si la suite est arithmétique ; si oui, donner sa raison.
\item Pour chaque suite arithmétique, interpréter la raison dans le contexte.
\end{enumerate}
}{%
\begin{enumerate}[label=\alph*.]
\item \textbf{Suite $(a_n)$.} On développe : $a_n=6n+21-n=5n+21$. Alors $a_{n+1}-a_n=5(n+1)+21-(5n+21)=5$. La différence est constante : $(a_n)$ est arithmétique de raison $5$ (premier terme $a_0=21$).

\textbf{Suite $(b_n)$.} $b_{n+1}-b_n=2(n+1)^2-(n+1)+4-(2n^2-n+4)=2(2n+1)-1=4n+1$. Cette différence dépend de $n$ : par exemple $b_0=4$, $b_1=5$, $b_2=10$, donc $b_1-b_0=1$ et $b_2-b_1=5$. La suite $(b_n)$ n'est pas arithmétique.

\textbf{Suite $(c_n)$.} $c_{n+1}-c_n=150-7(n+1)-(150-7n)=-7$. La suite $(c_n)$ est arithmétique de raison $-7$ (premier terme $c_0=150$).
\item Chaque semaine, l'entrepôt prépare $5$ colis de plus que la semaine précédente. Chaque semaine, le stock diminue de $7$ palettes (raison négative : diminution).
\end{enumerate}
}{4}

\mxexdef{2}{Terme général et termes éloignés (arithmétique)}{%
Une entreprise de livraison remplace progressivement ses véhicules diesel par des véhicules électriques. Au début de l'année 2025 (rang $0$), elle possède $186$ véhicules diesel et elle en retire $4$ chaque année. On note $d_n$ le nombre de véhicules diesel au début de l'année $2025+n$ ; la suite $(d_n)$ est arithmétique.
\begin{enumerate}[label=\alph*.]
\item Donner le premier terme $d_0$ et la raison de la suite $(d_n)$.
\item Exprimer $d_n$ en fonction de $n$ pour tout entier naturel $n$.
\item Calculer $d_{12}$ et $d_{30}$. Interpréter $d_{30}$ par une phrase.
\end{enumerate}
}{%
\begin{enumerate}[label=\alph*.]
\item $d_0=186$ ; on retire $4$ véhicules par an, donc $d_{n+1}=d_n-4$ : la raison est $r=-4$.
\item Pour une suite arithmétique de premier terme $d_0$ : $d_n=d_0+nr$, soit $d_n=186-4n$.
\item $d_{12}=186-4\times12=186-48=138$ et $d_{30}=186-4\times30=186-120=66$.

Le rang $30$ correspond à l'année $2025+30=2055$ : selon ce modèle, l'entreprise possèdera encore $66$ véhicules diesel au début de l'année $2055$.
\end{enumerate}
}{4}

\mxexdef{3}{Somme de termes consécutifs d'entiers}{%
\begin{enumerate}[label=\alph*.]
\item Calculer $\displaystyle\sum_{i=1}^{365}i=1+2+\cdots+365$.
\item Calculer $\displaystyle\sum_{k=100}^{365}k=100+101+\cdots+365$. Préciser le nombre de termes de cette somme.
\item Pour se constituer une épargne, Lova applique un défi : le $k$-ième jour de l'année, elle dépose $k$ centimes dans une tirelire ($1$ centime le 1\textsuperscript{er} jour, $2$ centimes le 2\textsuperscript{e} jour, etc.). Quelle somme, en euros, aura-t-elle déposée au total sur une année de $365$ jours ? Et entre le 100\textsuperscript{e} et le 365\textsuperscript{e} jour inclus ?
\end{enumerate}
}{%
\begin{enumerate}[label=\alph*.]
\item Les entiers de $1$ à $365$ forment une suite arithmétique de raison $1$ ; il y a $365$ termes. Donc
$\displaystyle\sum_{i=1}^{365}i=365\times\frac{1+365}{2}=365\times183=\num{66795}$.
\item De $100$ à $365$, il y a $365-100+1=266$ termes. Donc
$\displaystyle\sum_{k=100}^{365}k=266\times\frac{100+365}{2}=266\times\num{232,5}=\num{61845}$.

Contrôle : $\num{66795}-\num{61845}=\num{4950}=\dfrac{99\times100}{2}$, qui est bien la somme des entiers de $1$ à $99$.
\item Le total déposé sur l'année est $\num{66795}$ centimes, soit $\num{667,95}$~€. Entre le 100\textsuperscript{e} et le 365\textsuperscript{e} jour, Lova dépose $\num{61845}$ centimes, soit $\num{618,45}$~€.
\end{enumerate}
}{5}

\mxexdef{4}{Terme général d'une suite géométrique}{%
Une application est testée avant son lancement. La semaine $0$, $3$ testeurs l'installent. On suppose que, chaque semaine, le nombre de nouveaux testeurs est multiplié par $4$. On note $v_n$ le nombre de nouveaux testeurs de la semaine $n$ ; la suite $(v_n)$ est géométrique.
\begin{enumerate}[label=\alph*.]
\item Donner le premier terme et la raison de la suite $(v_n)$, puis exprimer $v_n$ en fonction de $n$.
\item Calculer $v_5$ et $v_9$.
\item Le modèle vous semble-t-il réaliste sur une longue durée ? Justifier par une phrase.
\end{enumerate}
}{%
\begin{enumerate}[label=\alph*.]
\item $v_0=3$ et $v_{n+1}=4v_n$ : la raison est $q=4$. Pour une suite géométrique, $v_n=v_0q^n$, donc $v_n=3\times4^n$.
\item $v_5=3\times4^5=3\times\num{1024}=\num{3072}$ et $v_9=3\times4^9=3\times\num{262144}=\num{786432}$.
\item Non : selon le modèle, il y aurait $\num{786432}$ nouveaux testeurs la semaine $9$, et plus de $3$ millions la semaine suivante ($v_{10}=\num{3145728}$). Une croissance aussi rapide ne peut pas durer : le nombre d'utilisateurs potentiels est limité.
\end{enumerate}
}{4}

\mxexdef{5}{Terme général et somme géométrique}{%
Une association de reboisement plante $\num{1500}$ arbres la première année de son programme. Grâce à de nouveaux bénévoles, le nombre d'arbres plantés double chaque année. On note $u_0=\num{1500}$ le nombre d'arbres plantés la première année et $u_n$ le nombre d'arbres plantés la $(n+1)$-ième année.
\begin{enumerate}[label=\alph*.]
\item Justifier que $(u_n)$ est une suite géométrique. Exprimer $u_n$ en fonction de $n$ et en déduire $u_6$.
\item Calculer $S=u_0+u_1+\cdots+u_6$ et interpréter ce résultat dans le contexte.
\end{enumerate}
}{%
\begin{enumerate}[label=\alph*.]
\item Doubler revient à multiplier par $2$ : $u_{n+1}=2u_n$. La suite $(u_n)$ est géométrique de premier terme $u_0=\num{1500}$ et de raison $q=2$. Donc $u_n=\num{1500}\times2^n$ et $u_6=\num{1500}\times2^6=\num{1500}\times64=\num{96000}$.
\item La somme comporte $7$ termes (de $u_0$ à $u_6$) et $q\neq1$ :
\[
S=u_0\,\frac{1-q^{7}}{1-q}=\num{1500}\times\frac{1-2^7}{1-2}=\num{1500}\times\frac{-127}{-1}=\num{190500}.
\]
Sur les sept premières années du programme, l'association a planté $\num{190500}$ arbres au total.
\end{enumerate}
}{5}

\mxexdef{6}{Choisir entre modèle arithmétique et géométrique}{%
Pour chaque proposition, dire si elle est vraie ou fausse. Justifier la réponse.
\begin{enumerate}
\item Le prix d'un abonnement annuel de transport augmente de $\num{1,50}$~€ chaque année. Cette situation peut être modélisée par une suite arithmétique de raison $\num{1,5}$.
\item Une voiture perd $\pc{15}$ de sa valeur chaque année. Cette situation peut être modélisée par une suite géométrique de raison $\num{1,15}$.
\item Une somme est placée sur un livret à intérêts composés au taux annuel de $\pc{3}$, sans aucun dépôt ni retrait. Le capital peut être modélisé par une suite géométrique de raison $\num{1,03}$.
\item Une cuve de récupération d'eau de pluie contient $\num{1000}$~L. Chaque jour, elle perd $\pc{5}$ de son contenu, puis on y ajoute $40$~L. Le volume peut être modélisé par une suite géométrique de raison $\num{0,95}$.
\item La population d'une commune augmente de $\pc{2}$ par an. Cette situation peut être modélisée par une suite arithmétique de raison $2$.
\item Un prêt de $\num{2000}$~€ est accordé à intérêts simples au taux annuel de $\pc{6}$. La somme due peut être modélisée par une suite arithmétique de raison $120$.
\item Le nombre de visites mensuelles d'un site internet diminue de $300$ chaque mois. Cette situation peut être modélisée par une suite géométrique.
\end{enumerate}
}{%
\begin{enumerate}
\item \textbf{Vrai.} On ajoute $\num{1,50}$~€ chaque année : $u_{n+1}=u_n+\num{1,5}$, la variation absolue est constante.
\item \textbf{Faux.} Une baisse de $\pc{15}$ correspond au coefficient $1-\dfrac{15}{100}=\num{0,85}$ : la suite est géométrique de raison $\num{0,85}$, et non $\num{1,15}$ (qui traduirait une hausse).
\item \textbf{Vrai.} Chaque année le capital est multiplié par $1+\dfrac{3}{100}=\num{1,03}$ : $C_{n+1}=\num{1,03}\,C_n$.
\item \textbf{Faux.} On a $V_{n+1}=\num{0,95}\,V_n+40$ : une opération additive s'ajoute à la multiplication. Avec $V_0=\num{1000}$ : $V_1=990$ et $V_2=\num{980,5}$ ; $\dfrac{V_1}{V_0}=\num{0,99}$ et $\dfrac{V_2}{V_1}\approx\num{0,9904}$ ne sont pas égaux, la suite n'est pas géométrique (elle n'est pas non plus arithmétique : $-10\neq-\num{9,5}$).
\item \textbf{Faux.} Une hausse de $\pc{2}$ par an est une évolution relative constante : la suite est géométrique de raison $\num{1,02}$.
\item \textbf{Vrai.} Aux intérêts simples, les intérêts sont calculés sur la somme initiale : $\num{2000}\times\num{0,06}=120$~€ s'ajoutent chaque année. La somme due augmente donc de $120$~€ par an.
\item \textbf{Faux.} La diminution est absolue et constante ($-300$ par mois) : la suite est arithmétique de raison $-300$.
\end{enumerate}
}{6}

% -----------------------------------------------------------------------------
% Chapitre 2 -- Fonctions exponentielles
% -----------------------------------------------------------------------------
\mxexdef{7}{Prolonger une suite géométrique en fonction}{%
Au 1\textsuperscript{er} janvier 2024, la puissance des panneaux solaires installés dans une région est de $240$ mégawatts (MW). On estime que cette puissance augmente de $\pc{12}$ par an. Pour tout entier naturel $n$, on note $u_n$ la puissance installée, en MW, au 1\textsuperscript{er} janvier de l'année $2024+n$ ; on admet que $(u_n)$ est une suite géométrique.
\begin{enumerate}
\item Donner la raison de la suite $(u_n)$ et exprimer $u_n$ en fonction de $n$.
\item On souhaite estimer la puissance installée à des dates quelconques. Proposer l'expression d'une fonction $f$ donnant la puissance installée $x$ années après le 1\textsuperscript{er} janvier 2024, $x$ étant un nombre réel positif.
\item Estimer, à $\num{0,1}$~MW près, la puissance installée au 1\textsuperscript{er} janvier 2026, au 1\textsuperscript{er} avril 2027 et au 1\textsuperscript{er} juillet 2029. On assimile un trimestre à $\num{0,25}$ an.
\end{enumerate}
}{%
\begin{enumerate}
\item Une hausse de $\pc{12}$ correspond au coefficient $1+\dfrac{12}{100}=\num{1,12}$. La raison est $q=\num{1,12}$ et $u_n=240\times\num{1,12}^n$.
\item On prolonge la suite aux réels positifs : $f(x)=240\times\num{1,12}^x$ (on a bien $f(n)=u_n$ pour tout entier $n$).
\item \begin{itemize}
\item 1\textsuperscript{er} janvier 2026 : $x=2$ et $f(2)=240\times\num{1,12}^2=\num{301,056}\approx\num{301,1}$~MW.
\item 1\textsuperscript{er} avril 2027 : $3$ ans et un trimestre, soit $x=\num{3,25}$ ; $f(\num{3,25})=240\times\num{1,12}^{\num{3,25}}\approx\num{346,9}$~MW.
\item 1\textsuperscript{er} juillet 2029 : $5$ ans et demi, soit $x=\num{5,5}$ ; $f(\num{5,5})=240\times\num{1,12}^{\num{5,5}}\approx\num{447,6}$~MW.
\end{itemize}
Selon ce modèle, la puissance installée serait d'environ $\num{301,1}$~MW, $\num{346,9}$~MW et $\num{447,6}$~MW à ces trois dates.
\end{enumerate}
}{5}

\mxexdef{8}{Propriétés algébriques des puissances}{%
Écrire chacun des nombres suivants sous la forme d'une seule puissance de $3$ :
\[
a=3^{-2}\times3^{6};\qquad b=\left(3^{\num{1,5}}\right)^4;\qquad
c=27\times9^{-2};\qquad d=\frac{3^8}{9^3}.
\]
}{%
\begin{itemize}
\item $a=3^{-2+6}=3^4$ (règle $a^xa^y=a^{x+y}$).
\item $b=3^{\num{1,5}\times4}=3^6$ (règle $(a^x)^y=a^{xy}$).
\item $27=3^3$ et $9=3^2$, donc $9^{-2}=\left(3^2\right)^{-2}=3^{-4}$ ; ainsi $c=3^3\times3^{-4}=3^{-1}$.
\item $9^3=\left(3^2\right)^3=3^6$, donc $d=\dfrac{3^8}{3^6}=3^{8-6}=3^2$.
\end{itemize}
On peut contrôler à la calculatrice, par exemple $c=\dfrac{27}{81}=\dfrac13=3^{-1}$ et $d=\dfrac{\num{6561}}{729}=9=3^2$.
}{3}

\mxexdef{9}{Sens de variation de $t\mapsto ka^t$}{%
Déterminer le sens de variation sur $\mathbb{R}$ des fonctions définies par :
\[
f:t\mapsto\num{3,5}\times\num{0,82}^t;\qquad
g:t\mapsto-\num{0,6}\times\num{1,15}^t;\qquad
h:t\mapsto-4\times\num{0,7}^t.
\]
}{%
\begin{itemize}
\item $0<\num{0,82}<1$, donc $t\mapsto\num{0,82}^t$ est strictement décroissante ; le coefficient $\num{3,5}$ est positif et conserve le sens de variation : $f$ est strictement décroissante.
\item $\num{1,15}>1$, donc $t\mapsto\num{1,15}^t$ est strictement croissante ; le coefficient $-\num{0,6}$ est négatif et inverse le sens de variation : $g$ est strictement décroissante.
\item $0<\num{0,7}<1$, donc $t\mapsto\num{0,7}^t$ est strictement décroissante ; le coefficient $-4$ est négatif et inverse le sens de variation : $h$ est strictement croissante.
\end{itemize}
Contrôle numérique : $h(0)=-4$, $h(1)=-\num{2,8}$ et $h(2)=-\num{1,96}$ : les valeurs augmentent.
}{3}

\mxexdef{10}{Taux d'évolution annuel moyen}{%
En cinq ans, une usine a réduit sa consommation d'eau de $\pc{27}$.
\begin{enumerate}[label=\alph*.]
\item Donner le coefficient multiplicateur global sur les cinq ans.
\item Calculer le taux d'évolution annuel moyen de la consommation d'eau, à $\pc{0,1}$ près.
\item Vérifier le résultat en appliquant cinq fois ce coefficient moyen arrondi.
\end{enumerate}
}{%
\begin{enumerate}[label=\alph*.]
\item Une baisse de $\pc{27}$ correspond au coefficient global $Q=1-\dfrac{27}{100}=\num{0,73}$.
\item Si $q$ est le coefficient annuel moyen, alors $q^5=\num{0,73}$, donc $q=\num{0,73}^{\frac15}\approx\num{0,9390}$. Le taux annuel moyen est $t=q-1\approx-\num{0,061}$ : la consommation d'eau a baissé en moyenne d'environ $\pc{6,1}$ par an.
\item $\num{0,939}^5\approx\num{0,7300}$, ce qui correspond bien à une baisse globale d'environ $\pc{27}$.
\end{enumerate}
}{4}

\mxexdef{11}{Lectures graphiques sur des courbes $a^x$}{%
On a représenté ci-dessous, dans un repère, les fonctions $f$ et $g$ définies sur $\mathbb{R}$ par $f(x)=\num{1,8}^x$ et $g(x)=\num{0,45}^x$, ainsi que la droite d'équation $y=3$.
\begin{center}
\begin{tikzpicture}
\begin{axis}[width=.78\textwidth,height=6.5cm,axis lines=middle,xmin=-3,xmax=3,ymin=0,ymax=6,
  xtick={-3,...,3},ytick={0,...,6},minor x tick num=4,minor y tick num=1,
  grid=both,major grid style={black!25},minor grid style={black!10},xlabel={$x$},ylabel={$y$}]
\addplot[mxteacher,very thick,samples=120,domain=-3:3]{0.45^x};
\addplot[mxwarn,very thick,samples=120,domain=-3:3]{1.8^x};
\addplot[mxmint,thick,dashed,domain=-3:3]{3};
\node[mxteacher,font=\small] at (axis cs:-1.45,4.8) {$\mathcal{C}_1$};
\node[mxwarn,font=\small] at (axis cs:2.2,5.2) {$\mathcal{C}_2$};
\node[mxmint,font=\small,fill=white,inner sep=1pt] at (axis cs:0.5,3.3) {$y=3$};
\end{axis}
\end{tikzpicture}
\end{center}
\begin{enumerate}
\item Quelle courbe représente la fonction $f$ ? Justifier ce choix.
\item Résoudre graphiquement l'équation $f(x)=3$. À l'aide d'un tableau de valeurs de la calculatrice, donner une valeur approchée de la solution à $\num{0,01}$ près.
\item Déterminer graphiquement l'ensemble des réels $x$ tels que $f(x)\geqslant3$, puis l'ensemble des réels $x$ tels que $g(x)\geqslant3$.
\end{enumerate}
}{%
\begin{enumerate}
\item $\num{1,8}>1$, donc $f$ est strictement croissante ; $0<\num{0,45}<1$, donc $g$ est strictement décroissante. La courbe croissante $\mathcal{C}_2$ représente $f$ (et $\mathcal{C}_1$ représente $g$). Les deux courbes passent par le point $\pt{0;1}$.
\item La courbe $\mathcal{C}_2$ coupe la droite $y=3$ en un point d'abscisse lue proche de $\num{1,9}$. Avec la calculatrice : $\num{1,8}^{\num{1,86}}\approx\num{2,984}$ et $\num{1,8}^{\num{1,87}}\approx\num{3,002}$, donc la solution est comprise entre $\num{1,86}$ et $\num{1,87}$ ; la valeur $\num{1,87}$ donne l'image la plus proche de $3$. La solution est $x\approx\num{1,87}$.
\item $f(x)\geqslant3$ lorsque $\mathcal{C}_2$ est au-dessus de la droite $y=3$, c'est-à-dire pour $x\geqslant\num{1,87}$ environ : $\intfo{\num{1,87}}{+\infty}$ (valeur approchée de la borne).

$g$ est décroissante et $\mathcal{C}_1$ coupe la droite $y=3$ pour une abscisse proche de $-\num{1,4}$ (calculatrice : $\num{0,45}^{-\num{1,38}}\approx\num{3,010}$ et $\num{0,45}^{-\num{1,37}}\approx\num{2,986}$, soit $x\approx-\num{1,38}$). Ainsi $g(x)\geqslant3$ pour $x\leqslant-\num{1,38}$ environ : $\intof{-\infty}{-\num{1,38}}$ (valeur approchée de la borne).
\end{enumerate}
}{5}

\mxexdef{12}{Inéquations avec une même base}{%
Résoudre sur $\mathbb{R}$ les inéquations suivantes.
\begin{enumerate}[label=\alph*.]
\item $\num{0,7}^x\leqslant\num{0,7}^{-1}$
\item $\num{2,5}^{x+1}>\num{2,5}^4$
\item $\num{0,5}^x\geqslant\num{0,25}$
\end{enumerate}
}{%
\begin{enumerate}[label=\alph*.]
\item $0<\num{0,7}<1$ : la fonction $x\mapsto\num{0,7}^x$ est strictement décroissante, elle inverse l'ordre. Donc $\num{0,7}^x\leqslant\num{0,7}^{-1}\iff x\geqslant-1$. Ensemble des solutions : $\intfo{-1}{+\infty}$.
\item $\num{2,5}>1$ : la fonction $x\mapsto\num{2,5}^x$ est strictement croissante, elle conserve l'ordre. Donc $\num{2,5}^{x+1}>\num{2,5}^4\iff x+1>4\iff x>3$. Ensemble des solutions : $\intoo{3}{+\infty}$.
\item $\num{0,25}=\num{0,5}^2$, donc l'inéquation s'écrit $\num{0,5}^x\geqslant\num{0,5}^2$. Comme $x\mapsto\num{0,5}^x$ est strictement décroissante : $x\leqslant2$. Ensemble des solutions : $\intof{-\infty}{2}$.
\end{enumerate}
}{3}

\mxexdef{13}{Modéliser par $f(t)=ka^t$ et calculer}{%
Juste après une injection, la concentration d'un médicament dans le sang d'un patient est de $80$~mg/L. Cette concentration diminue ensuite de $\pc{15}$ chaque heure. Il s'agit d'une décroissance exponentielle.

On la modélise par une fonction $C$ du temps $t$, exprimé en heures, vérifiant $C(t)=ka^t$.
\begin{enumerate}
\item Déterminer les constantes $k$ et $a$.
\item Calculer la concentration au bout de $3$~h~$30$~min, à $\num{0,1}$~mg/L près.
\item Le médecin prévoit une nouvelle injection dès que la concentration passe sous $50$~mg/L. Est-ce déjà le cas au bout de $3$~h~$30$~min ? Répondre par une phrase.
\end{enumerate}
}{%
\begin{enumerate}
\item $C(0)=k\times a^0=k$, et $C(0)=80$, donc $k=80$. Une baisse de $\pc{15}$ par heure correspond au coefficient $1-\dfrac{15}{100}=\num{0,85}$, donc $a=\num{0,85}$ et $C(t)=80\times\num{0,85}^t$.
\item $3$~h~$30$~min $=\num{3,5}$~h. $C(\num{3,5})=80\times\num{0,85}^{\num{3,5}}\approx\num{45,3}$~mg/L.
\item $\num{45,3}<50$ : au bout de $3$~h~$30$~min, la concentration est déjà passée sous $50$~mg/L, la nouvelle injection est donc prévue (d'ailleurs $C(3)\approx\num{49,1}$~mg/L, la concentration est sous ce seuil dès $3$~h).
\end{enumerate}
}{4}

\mxexdef{14}{Taux moyen de deux évolutions successives}{%
\begin{enumerate}
\item Le cours d'une monnaie numérique augmente de $\pc{125}$ une année, puis baisse de $\pc{36}$ l'année suivante. Quel est le taux d'évolution annuel moyen sur ces deux années ?
\item Le prix d'une matière première augmente de $\pc{80}$ un mois, puis baisse de $\pc{55}$ le mois suivant. Quel est le taux d'évolution mensuel moyen ?
\item Un élève affirme que, dans la question 2, le taux moyen est $\dfrac{80-55}{2}=\pc{12,5}$. Expliquer par une phrase pourquoi il se trompe.
\end{enumerate}
}{%
\begin{enumerate}
\item Les coefficients sont $1+\dfrac{125}{100}=\num{2,25}$ et $1-\dfrac{36}{100}=\num{0,64}$. Le coefficient global est $\num{2,25}\times\num{0,64}=\num{1,44}$. Le coefficient annuel moyen $q$ vérifie $q^2=\num{1,44}$, donc $q=\sqrt{\num{1,44}}=\num{1,2}$ : hausse annuelle moyenne de $\pc{20}$.
\item Les coefficients sont $\num{1,8}$ et $\num{0,45}$ ; le coefficient global est $\num{1,8}\times\num{0,45}=\num{0,81}$. Le coefficient moyen est $\sqrt{\num{0,81}}=\num{0,9}$, soit une baisse mensuelle moyenne de $\pc{10}$.
\item Les évolutions successives se composent en multipliant les coefficients, pas en additionnant les taux : la baisse de $\pc{55}$ s'applique à un prix déjà augmenté. Sur les deux mois, le prix est multiplié par $\num{0,81}$ (baisse globale de $\pc{19}$), ce qui correspond à une baisse moyenne de $\pc{10}$ par mois et non à une hausse de $\pc{12,5}$.
\end{enumerate}
}{5}
